% !TEX root =  ../Dissertation.tex

\chapter{Conclusion and Future Work}
\label{ch:conclusion}

I built a portfolio system that is trained on the quality of its own decisions rather than on forecast accuracy like with MSE. A differentiable ledger whose realized excess Sharpe is the loss, a frozen Simple Cycle Reservoir to map to a higher feature dimension, and a gated mixture of experts whose regimes are discovered by that loss itself. To test this fairly I build various helper functions like a virtual ledger that takes portfolio suggestions and run them through a fixed path for that given time range, accounting for transaction cost, trading days etc. All models including the baselines, our proposed models, and the LSTM run through this ledger. For reproducibility I set fixed seeds, various deterministic algorithms, and to verify an append-only run registry. Finally a holdout fire walled from the first day and only touched once.

Three central findings are noticed. First, regimes discovered by the decision objective are completely different from statistically defined ones, and we can even say why: the statistical gates track risk, ours tracks recent market direction, and the symmetry argument of Chapter~\ref{ch:discussion} shows given the zero-bias linear gate and odd reservoir state used here, the architecture cannot represent a purely even, sign-invariant risk partition from the reservoir state alone. Second, the difference does not pay. No learned arm beats(as in is significantly better) the best naive baseline on either universe, no difference involving a learned arm is significant, and the holdout reproduces the same out of sample. Third, hyperparameter selection does not transfer on this problem meaning for either encoder, configurations chosen on validation data come back no better, and sometimes worse, than untuned defaults. I regard the honest null, delivered with its mechanism, as a result rather than a disappointment, and the evaluation protocol that made it trustworthy as a contribution in its own right.

\section{Future work}
\label{sec:concl-future}
First, as promised in Section~\ref{sec:lim-evaluation}: replace the single-path walk-forward with combinatorial purged cross-validation. It produces a distribution of backtest paths and, from it, an estimate of the probability that a selected strategy is overfit. Since our central methodological finding is that selection does not transfer, the natural next study measures that probability rather than inferring it from held-out checks.

The linear gate works like an odd function and cannot express a risk partition (Eq.~\ref{eq:odd-state}). The open question is what a gate free of that constraint discovers. Two ways to do this is to either use a small nonlinear gate, or to use simple even features of the state such as its squares and then re-run through the same $\eta^2$ analysis. If the nonlinear gate still organises on direction, the revealed preference of the decision objective is confirmed. If it switches to risk, the architecture, not the objective, was the binding constraint.

The single reservoir reads a fixed-width vector, which is why the universe must be frozen at inception. This was the reason why I chose these two universes to drive the encoder, and why a changing index like the Dow was out of reach (Section~\ref{sec:universes}). A natural redesign gives each asset its own small cycle reservoir, with shared or individual hyperparameters, and lets the heads read the pooled per-asset states. An asset's reservoir then simply starts when the asset comes into the market. Assets can thus freely enter and leave, the input dimension is no longer static, and the method extends to evolving universes.

We see a large gap in the training and validation sharpe values at all K values. This is likely due to drift, the optimal distribution/regime structure , etc changes from one year to next, which is what the gap shows. One solution would be to have an objective that rewards transfer rather than fit, such as maximising the \emph{worst} sub-period Sharpe of the training span instead of the pooled one, so that a configuration is preferred only if it survives every regime it has already seen.